\documentclass[a4paper]{article}
\usepackage{amsmath}
\usepackage{amssymb}
\usepackage{geometry}
\usepackage{hyperref}

\geometry{margin=25mm}

\title{Electron Orbits from the Signed Electrostatic Force:\\
Covalent Bonds, Metallic Bonds, and Crystals}

\author{https://github.com/hypernumbernet}
\date{\today}

\begin{document}
\maketitle

\begin{abstract}
On the equal-scale locus the electromagnetic force of double spacetime theory is the Coulomb force divided by the interference factor \(\gamma_s\). The sign of \(\gamma_s\) is the repulsion. An unlike pair attracts where \(\gamma_s>0\) and repels where \(\gamma_s<0\); a like pair does the opposite. No centrifugal term is added to cancel an attraction. A single electron--proton pair carries both channels. Where the two interference factors have the same sign, the gravitational term of the two masses and the electric term point the same way, and the radial force does not vanish. At the Bohr radius both factors are positive, and the mass term is smaller than the charge term by a factor strictly between \(4.40\times 10^{-40}\) and \(4.42\times 10^{-40}\). In the outer well a circular path is radially restoring above a radius that lies strictly between \(5\ell/6\) and \(\ell\), and unstable between that radius and the first node. The circular paths of that well have their least energy at the same radius, and the value is negative. One electron about a nucleus of any positive charge keeps that radius, and the energy of the path is the charge times the one-proton value. Two electrons on a diameter, four at the vertices of a tetrahedron, six at the vertices of an octahedron, and eight at the vertices of a cube, each have exactly one radial balance in the first repulsive shell; the diameter is an angular maximum, and the counts \(2\), \(8\), \(18\), and \(32\) are a dual spinor and its spin-reversed partner on the modes through degrees \(0\), \(1\), \(2\), and \(3\), not indices of the nodes. Degree by degree the seats number \(2\), \(6\), \(10\), and \(14\). On each of these patterns the square of the angular momentum, the response of the pattern when the direction is turned, is the product of the degree and one more than the degree, so the four degrees carry \(0\), \(2\), \(6\), and \(12\). A seat is a pattern on directions. What is called an electron orbital is an electron configuration, not an orbit. An electron may still travel, and in a covalent bond the places are fixed. Only the monopole is independent of direction, and every higher seat is positive in some directions and negative in others. The restoring radii of the diameter, of the tetrahedron, of the octahedron, and of the cube lie where \(\gamma_s<0\). In the outer well the same two electrons are not simultaneously force-free once the nuclear charge is at least \(1/4\), in any placement: on opposite rays the repulsion is strictly less than a quarter of the attraction of unit nuclear charge at every finite radius. Where the interference factor of a separation is positive, the antipode is an angular minimum, so the electrons spread and the radius is not restored. The shell of two, where the nucleus attracts, is the two monopole seats, not a zero of the signed force. In the outer well, where the nucleus attracts, a tetrahedron, an octahedron, a cube, and an icosahedron have no radial root once the nuclear charge is at least \(1\), \(2\), \(3\), and \(5\): every chord is longer than the radius, and the outward coefficient lies below that charge. A regular dodecahedron has a chord shorter than the radius, and the far-field coefficient of its twenty vertices lies strictly between \(15/2\) and \(8\). Before that chord meets the first node the response stays strictly above \(7\): a charge of at most \(7\) is pushed outward on that whole interval, while for every \(Z\ge 8\) the response equals \(Z\) at exactly one such phase, and the radius there lies strictly between \(7\ell/10\) and \(10\ell/7\). Twenty is the number of vertices, not one of the four capacities. On a segment longer than twice the first-node radius, an electron in the outer well with respect to both nuclei is drawn from the midpoint toward the larger charge. The one cancellation lies closer to the smaller charge and is not restored. Two protons and two electrons do, and in the first repulsive shell they do so at exactly one square, which restores the diagonal. While the proton separation and the electron separation both lie in the outer well, the two forces do not vanish together at any ratio of those separations. While the electron--proton separation lies in the outer well, the two forces do not vanish together. Helium, a nucleus of charge two with two electrons, remains the shell of two: the same family with charge two on each nucleus has no equilibrium in that well, and the square is not an equilibrium wherever both interference factors are nonzero. The water molecule places one bonding electron, closer than its proton, on each of two rays from the nucleus: while those electrons repel, the tangential force on each stays toward the wider angle so long as every further electron stands on the ray or repels it from the side toward the other bond, and the straight line restores those electrons. On the straight line, with each electron between the nucleus and its proton, the outward sum is negative once the nuclear charge is at least \(1/4\). When the electron lies at least a quarter as far from the nucleus as the proton and the separation at which the protons are as far apart as each is from the other electron still lies in the outer well, the tangential force on each proton vanishes at an angle strictly between that equal separation and the straight line. The angle between the rays at that equal separation is strictly less than a third of a straight angle. Four shared pairs on the rays of a regular tetrahedron, the electrons at one distance and the protons at one greater distance, exert no tangential force. That angle has cosine \(-1/3\) and is greater than a right angle. A non-shared pair has no proton on its ray, so its electron is driven inward. While that electron stays at least \(1/\sqrt{2}\) times as far from the nucleus as a neighbor and the separations lie in the outer well, it pushes the bonds aside more strongly than a shared pair at the bonding distance, and two such pairs contribute twice the narrowing of one. Protons sit on one diagonal and electrons on the other, and the electrostatic potential has a positive-definite Hessian there. Like charges attract along the diagonals and unlike charges repel along the sides. The diagonal is \(\sqrt{2}\) times the side, the diagonal force is \(4.3532\,k e^{2}/\ell^{2}\), and the stiffness of the diagonal is \(8.178\,k e^{2}/\ell^{3}\). A nearest-neighbor bond that carries only the fraction \(2/z\) of that pair has further points where every electrostatic component vanishes. For eight and for twelve neighbors those points are saddles, and the partial forces stand in the ratios \(0.04323\) and \(0.1704\) to the covalent diagonal. In the outer well of every ion the ions alone hold no electron. Displaced along three perpendicular directions, the electron feels along each a change of force, and the three rates at which the displacement is pushed further sum to a positive value, which the Coulomb law makes zero. Further electrons do not reverse this so long as each lies in the outer well and farther away than a distinct ion. At the center of a regular octahedron, a cube, or a cuboctahedron of ions the force on an electron vanishes and every displacement is pushed further, while the same shells restore every displacement of a positive ion. The ions are held in place, and the electrons are free. A cubical truncation of a simple, body-centered, or face-centered lattice of like ions holds the ion in every direction and pushes an electron off, to any number of shells, while every phase lies below \(\pi\) and off the first node. The rate has the sign of \(\sin x\), so it reverses at radius \(\ell/(2\pi)\). The nearest two shells of rock salt, cesium chloride, zincblende, and both sites of fluorite push a positive ion off throughout the outer well for the ordinary ionic charges. A rotation through \(2\pi/n\) preserves a three-dimensional lattice only for \(n\in\{1,2,3,4,6\}\), so a crystal site has no fivefold axis. The side of that square is not the Bohr radius. Placing the side at that radius would put the first node strictly between \(91/40\) and \(12/5\) Bohr radii. The laboratory sign at the Bohr radius is attraction, so that radius lies in the outer well. The square lies in the reversed shell. The same layered sign, read in the gravitational channel, is the exclusion of a shared seat: two electron masses placed inside the first gravitational node meet a simple pole. The smallness of \(G m_e^{2}\) thins that neighborhood and does not remove the pole. A restoring radius is a sign boundary of the outward force, and the electron oscillates across it. On the square the ratio of the electron wavelength to the diagonal wavelength is \(45.08\), the square root of the ratio of \(K/m\) on the electron coordinate to \(K/m\) on the diagonal. The ratio does not use an identification of the side. A transition frequency is an eigenvalue difference and cannot scale with temperature. A stationary orbital fixes the electron's motion, so the velocity of the electron relative to the proton is not a quantity the orbital computes, and that difference does not change with the temperature. On a harmonic ladder the mean number of quanta above the ground state is \(1/(e^{x}-1)\), with \(x\) the spacing in units of \(k_B T\). In a large enclosure the number of modes grows as the square of the frequency. That growth is the geometry of the enclosure. The wavelength density then has a single maximum between \(9/2\) and \(5\) in units of \(k_B T\). The laboratory value \(4.96\,k_B T\) is that maximum, so the peak wavelength times the temperature is constant. The integrated frequency density scales as \(T^{4}\), and the power of one classical root scales as \(T\). The constants \(h\), \(k_B\), and \(c\), and the growth of the mode count, are inputs.
\end{abstract}

\section{Introduction}

The electron-shell section of double spacetime theory writes an electrostatic force whose denominator is the scalar interference factor \(\gamma_s\)~\cite{dst-main}. At large separation the factor tends to one and the Coulomb law is recovered. Closer in, \(\gamma_s\) passes through zero on a discrete set of radii and is negative between those nodes. A node is a divergence of the force. The radial potential on the outer side of a node is unbounded below.

A purely attractive Coulomb force does not vanish at any finite radius. The classical circular orbit cancels it by a centrifugal term \(L^{2}/(m r^{3})\), the kinematics of a conserved angular momentum. That term is not added here. Section~\ref{sec:nocent} places the change of a circular path in the outer well, from radially restoring to unstable, at a radius strictly between \(5\ell/6\) and \(\ell\). That radius carries the least energy of those circular paths, and the energy is negative. For one electron about a nucleus of any positive charge the radius is the same, and the energy is the charge times the one-proton value. The outward contribution is already in the signed force. Where \(\gamma_s\) is negative an unlike pair repels, and where \(\gamma_s\) is positive a like pair repels. A cancellation of these contributions is a radial root when the charges and the geometry supply one. The principal capacities are the seats of Section~\ref{sec:shells}, and the shell of two is not such a root. Neither is a balance between Coulomb attraction and a centrifugal barrier.

The note computes that locus for a two-center arrangement of two protons and two electrons. In the outer well the arrangement does not balance at any ratio of the separations. While the electron--proton separation lies in the outer well, the two forces do not vanish together. The one square in the first repulsive shell restores the diagonal. Helium, a nucleus of charge two with two electrons, remains the shell of two: the same family with charge two on each nucleus has no equilibrium in that well, and the square is not an equilibrium wherever both interference factors are nonzero. Section~\ref{sec:water} places one bonding electron, closer than its proton, on each of two rays from the nucleus. While those electrons repel, the tangential force stays toward the wider angle so long as every further electron stands on the ray or repels it from the side toward the other bond, and the straight line restores them. On the straight line the outward sum is negative once the nuclear charge is at least \(1/4\). When the electron lies at least a quarter as far from the nucleus as the proton and the equal proton separation still lies in the outer well, the tangential force on each proton vanishes strictly between that separation and the straight line. The angle between the rays at the equal separation is strictly less than a third of a straight angle. Section~\ref{sec:angle} places four pairs on the rays of a regular tetrahedron. Their angle has cosine \(-1/3\) and is greater than a right angle, and four shared pairs, electrons at one distance and protons at one greater distance, exert no tangential force. A non-shared pair has no proton on its ray, and its electron is driven inward. While that electron remains at least \(1/\sqrt{2}\) times as far from the nucleus as a neighbor, it pushes the remaining bonds aside more strongly than a shared pair, and two such electrons contribute twice the narrowing of one. The same calculation covers a nearest-neighbor bond that carries only \(2/z\) of an electron pair. Section~\ref{sec:free} shows that in the outer well of every ion the ions alone hold no electron, whereas a symmetric shell of like ions holds an ion at its center. The electrons of a metal are free. Section~\ref{sec:crystal} extends that holding to every shell of a cubic truncation, and the nearest two shells of the common ionic crystals push the ion off. A further stability belongs to the gravitational channel. The main text reads the Pauli principle as the first repulsive layer of that channel, generated when two fermions overlap, not as a separate statistical postulate. Applied to electrons, the same layer expels two masses from one orbital seat. The electron mass is small, and the force that expels them is a pole. The equal-scale identification \(\alpha=\beta=\ell/r\) is the reading already used for the nuclear and Coulomb layers. The electromagnetic length is the reduced Compton wavelength of the electron. Bohr's \(n^{2}\) rule and the fine-structure constant are not derived. The same refusal covers the shell counts: \(2\), \(8\), \(18\), and \(32\) are a dual spinor and its spin-reversed partner on the monopole, the dipole, the quadrupole, and the octupole. The square of the angular momentum on those patterns, their response when the direction is turned, is \(0\), \(2\), \(6\), and \(12\). Only the monopole is the same in every direction. Every higher seat is positive in some directions and negative in others. What is called an electron orbital is an electron configuration, not an orbit. An electron may still travel, and in a covalent bond the places are fixed. Quantum mechanics already reads the orbital as a stationary pattern, and the signed force carries that reading through to the places. A one-center diameter, tetrahedron, octahedron, or cube closes only its radius, and only in the reversed layer. Two electrons in the outer well are not simultaneously force-free once the nuclear charge is at least \(1/4\). The laboratory well, in which the nucleus attracts, leaves a tetrahedron, an octahedron, a cube, and an icosahedron without a radial root once the nuclear charge is at least \(1\), \(2\), \(3\), and \(5\). A dodecahedron of twenty vertices has exactly one radial root in the outer well once the nuclear charge is at least \(8\), at a radius strictly between \(7\ell/10\) and \(10\ell/7\). Twenty is not one of the four capacities. The far-field coefficient of those vertices lies strictly between \(15/2\) and \(8\), and before the nearest chord meets the first node the response stays strictly above \(7\), so a charge of at most \(7\) is pushed outward on that whole interval. Section~\ref{sec:eneg} places an electron on the segment between two nuclei in that well. From the midpoint it is drawn toward the larger charge, and the one cancellation lies closer to the smaller charge. The restoring radius is also a sign boundary. Section~\ref{sec:heat} reads the radial oscillation across it. The frequency is \(\omega^{2}=K k e^{2}/(m\ell^{3})\): one signed balance, and a separate inertia for each coordinate that moves. A transition frequency is an eigenvalue difference and contains no temperature. A stationary orbital fixes the electron's motion, so the velocity of the electron relative to the proton is not a quantity the orbital computes, and that difference does not change with the temperature. The continuous spectrum of an enclosure is an occupation on a ladder of fixed spacing. Where a laboratory number enters, it is marked as an input.

\section{The signed force}
\label{sec:force}

On one axis the interference factor of the dual rotor is
\[
\gamma_s
  =\cosh\frac{\alpha}{2}\cos\frac{\beta}{2}
  -\sinh\frac{\alpha}{2}\sin\frac{\beta}{2}.
\]
The equal-scale locus sets \(\alpha=\beta=\ell/r\). In the dimensionless variable \(x=\ell/(2r)\),
\begin{equation}
\label{eq:gamma}
\gamma_s(x)=\cosh x\cos x-\sinh x\sin x,
\qquad
\gamma_s'(x)=-2\cosh x\sin x.
\end{equation}
Thus \(\gamma_s(0)=1\), and on each interval \([n\pi,(n+1)\pi]\) there is one simple zero. The first zero \(x_1\) is the unique root of
\begin{equation}
\label{eq:node}
\tan x=\coth x
\end{equation}
in \((\pi/4,\pi/2)\). Numerically \(x_1=0.93755\), so the outermost node lies at
\[
\frac{r_1}{\ell}=\frac{1}{2x_1}=0.53330.
\]
The second zero lies at \(x_2=3.9274\), and the inner edge of the first repulsive shell is \(r_2/\ell=0.12731\). Ordered by decreasing radius, the outer well \(r>r_1\) has \(\gamma_s>0\), the shell \(r_2<r<r_1\) has \(\gamma_s<0\), and the sign continues to alternate. The extremum of \(\gamma_s\) inside the first shell is \(\gamma_s(\pi)=-\cosh\pi\), at \(r=\ell/(2\pi)\). The extremum of the force lies at a larger radius (Section~\ref{sec:nocent}).

The force that increases the separation of two charges is
\begin{equation}
\label{eq:force}
F_r=\frac{k q_1 q_2}{\gamma_s(r)\,r^{2}},
\qquad
k=\frac{1}{4\pi\epsilon_0}.
\end{equation}
Like charges repel where \(\gamma_s>0\) and attract where \(\gamma_s<0\). Unlike charges do the opposite. Both directions are the same formula. Nothing further is added to produce a repulsion.

An electron and a proton keep the gravitational term of their masses beside this electric term. With \(\gamma_{s,G}\) the interference factor of the mass channel at \(\lambda_\pi\), the charged-pion Compton wavelength, the force that increases their separation is
\begin{equation}
\label{eq:both}
F_r
  =-\frac{G m_e m_p}{\gamma_{s,G}(r)\,r^{2}}
  -\frac{k e^{2}}{\gamma_s(r)\,r^{2}}.
\end{equation}
Where the two factors have the same sign, the two terms point the same way and the sum does not vanish. The Bohr radius lies on the outer plateau of both lengths. Both factors are positive there, both terms attract, and the mass term is smaller than the charge term by a factor strictly between \(4.40\times 10^{-40}\) and \(4.42\times 10^{-40}\).

The Taylor expansion at the origin follows from \(\gamma_s''(0)=-2\) and \(\gamma_s''''(0)=-4\),
\begin{equation}
\label{eq:series}
\gamma_s(x)=1-x^{2}-\frac16 x^{4}+O(x^{6}).
\end{equation}
Hence
\[
\frac{1}{\gamma_s(r)\,r^{2}}
  =r^{-2}+\frac{\ell^{2}}{4}\,r^{-4}+\frac{7\ell^{4}}{96}\,r^{-6}+O(r^{-8}).
\]
The correction to the Coulomb force begins at order \((\ell/r)^{2}\).

For \(r>r_1\) the pair potential that vanishes at infinity is the integral of~\eqref{eq:force}. The substitution \(s=\ell/(2x)\) gives \(ds/s^{2}=-(2/\ell)\,dx\), and therefore
\begin{equation}
\label{eq:potential}
\int_r^{\infty}\frac{ds}{\gamma_s(s)\,s^{2}}
  =\frac{2}{\ell}\int_0^{\ell/(2r)}\frac{dx}{\gamma_s(x)},
\end{equation}
provided the path stays in the outer well. The integral diverges as the upper limit approaches \(x_1\). Inside a shell the same formula does not continue: the path from infinity would cross a pole. A local force balance inside the shell is still well defined, because~\eqref{eq:force} requires only \(\gamma_s\neq 0\).

\section{No centrifugal term}
\label{sec:nocent}

Consider one proton and one electron. The radial force is
\[
F_r=-\frac{k e^{2}}{\gamma_s(r)\,r^{2}}.
\]
Wherever \(\gamma_s\) is finite and nonzero, \(F_r\) is nonzero. In the outer well it is attractive at every finite radius. In the first repulsive shell it is repulsive at every finite radius. There is no radius at which the electrostatic force vanishes, and none is introduced by hand.

The centripetal rearrangement \(v^{2}=k e^{2}/(m\gamma_s r)\) of the main paper records what a conserved angular momentum contributes in the outer well. It is a kinematic identity for a circular Newtonian orbit, and it is not an extra force in~\eqref{eq:force}. Where \(\gamma_s>0\), the inward force supplies that acceleration.

Counted with attraction positive, the dimensionless force is \(Y=4x^{2}/\gamma_s(x)\). In the outer well it is positive and strictly increases with the phase, so the attraction strengthens as the radius falls. The boundary of the well is the node. Inside the first repulsive shell the same force is negative. It diverges at either node and has exactly one extremum, at a phase strictly between \(5/2\) and \(\pi\). The radius there lies strictly between \(\ell/(2\pi)\) and \(\ell/5\), and the outward magnitude is strictly less than \(3\,k e^{2}/\ell^{2}\). At \(x=\pi\), where \(\gamma_s=-\cosh\pi\),
\begin{equation}
\label{eq:plateau}
\bigl|F\bigr|
  =\frac{4\pi^{2}}{\cosh\pi}\,\frac{k e^{2}}{\ell^{2}}
  =3.4057\,\frac{k e^{2}}{\ell^{2}},
\end{equation}
directed outward. The growth of \(1/r^{2}\) places the weaker extremum at a larger radius than this point. With no angular momentum, an unlike pair is driven toward the first node from both sides: the outer well pulls inward, and the first shell pushes outward. A like pair is driven off that node, and the wall of a shell in the laboratory well is that repulsion. A one-center electron is not parked on a force-free circle. A force-free locus appears only when several charges are present and their signed contributions cancel. Between two nuclei in the outer well that cancellation is not a seat (Section~\ref{sec:eneg}).

The circular path is radially restoring while \(d\ln Y/d\ln x<3\). On the outer well that inequality holds exactly when \(\gamma_s(x)>2x\cosh x\sin x\), and the two sides meet at exactly one phase, strictly between \(1/2\) and \(3/5\). The radius there lies strictly between \(5\ell/6\) and \(\ell\), outside the first node. Between that radius and the node the circular path is unstable.

The energy of a circular path adds the kinetic term of the centripetal identity to the outer-well potential~\eqref{eq:potential}. In units of \(k e^{2}/\ell\) it is \(x/\gamma_s(x)\) minus twice the integral of \(1/\gamma_s\) from zero to the phase. The energy falls as the phase rises through the restoring side, and it rises from the one phase already located toward the node. That phase carries the least circular energy in the outer well, and the value is negative. The force there remains attractive. The minimum selects a circular path, and the radial potential of a plunge stays unbounded below at the node.

The nuclear charge, counted in units of the proton charge, is the proton number \(Z\). One electron about that nucleus feels \(Z\) times the one-proton force. The centripetal identity and the outer-well potential each carry the same factor, so the circular energy, in units of \(k e^{2}/\ell\), is \(Z\) times the one-proton energy, and the value at that radius remains negative. The logarithmic derivative of the force is the one-proton derivative, so every positive charge restores the circular path on the same side of that radius. A larger charge makes this energy lower.

\section{Shell counts}
\label{sec:shells}

The integers \(2\), \(8\), \(18\), and \(32\) are not labels of the equal-scale nodes. Those nodes still march inward as \(\Theta(1/n)\). Bohr's ladder of radii \(n^{2}\) is not read off from them.

Two electrons at the ends of a diameter, about a nucleus of charge \(Z\ge 1\), have no radial balance in the outer well. In the first repulsive shell, where both the radius and the diameter lie between the first node and the second, there is exactly one, whichever \(Z\) is. There the nucleus pushes outward and the two electrons attract. On the side of larger radius the outward force points inward, and on the side of smaller radius it points outward, so the radius is restored. The angle is not. At a fixed radius the separation of the pair is greatest when the electrons are antipodal, and wherever the interference factor of that separation is negative an energy that increases with separation has a maximum there. A displacement that draws the two electrons toward one common seat is not restored by the electrostatic force.

The same absence is not special to the diameter. Two electrons whose distances to the nucleus and whose mutual separation all lie in the outer well are not simultaneously force-free once the nuclear charge is at least \(1/4\), whatever the placement. On opposite rays the repulsion is strictly less than a quarter of the attraction of unit nuclear charge at every finite radius, and every nucleus is at least that quarter. Where the interference factor of the separation is positive, the antipode is an angular minimum of an energy that falls as the separation grows, so the electrons spread. That spreading does not restore a radius. The closed shell of two, in the well where the nucleus attracts, is the two monopole seats. It is not a zero of the signed force.

Eight electrons at the vertices of a cube about the same nucleus have, again for every \(Z\ge 1\), exactly one radial balance in the window where the radius and the three separations---the edge, the face diagonal, and the body diagonal---all lie in the first repulsive shell. The outward force restores that radius.

Four electrons at the vertices of a regular tetrahedron about the same nucleus have, for every nuclear charge at least \(1\), exactly one radial balance in the window where the radius and the edge both lie in the first repulsive shell. Six electrons at the vertices of a regular octahedron have exactly one in the window where the radius, the edge, and the opposite vertex all lie there, again for every such charge. In each case the outward force restores the radius. Four and six are the numbers of vertices. They are not seat counts, and the six vertices are not the six dipole seats.

Eighteen is not the number of vertices of one cubic figure. Those figures have six, eight, or twelve vertices. Eighteen counts seats. A monopole is one pattern, a dipole three, a traceless quadrupole five, and a traceless octupole seven: one more than twice the degree. The square of the angular momentum on a pattern is the product of the degree and one more than the degree. It is the response of the pattern when the direction is turned, and on these four patterns it is \(0\), \(2\), \(6\), and \(12\). A dual spinor and its spin-reversed partner put one seat on each pattern. Taken degree by degree, those seats number \(2\), \(6\), \(10\), and \(14\). Through degree zero the count is \(2\), through degree one it is \(8\), through degree two it is \(18\), and through degree three it is \(32\). The same running total, through degree \(n-1\), is \(2n^{2}\) for \(n=1,2,3,4\). The four principal-shell capacities are these four numbers. The count does not locate a node, and no figure of eighteen or of thirty-two points is claimed here to be a minimum of the force.

A seat is a pattern on directions. The square is the answer of that pattern to a turn of the direction. The monopole takes the same value in every direction, and the square is \(0\). A dipole changes sign across the plane through the nucleus perpendicular to its axis: one side is positive, the other is negative, and that plane is the set of directions on which the value is zero; the square is \(2\). A quadrupole keeps its value when the direction is reversed. Where it is not zero, it is positive in some directions and negative in others, and the value is zero along some direction in their plane; the square is \(6\). An octupole changes sign when the direction is reversed: the value in the opposite direction is the negative of the value, and the square is \(12\). The mode fixes the pattern.

What is called an electron orbital is an electron configuration. The seat places the electron, and the placement is not an orbit. An electron may still travel. The circular path of Section~\ref{sec:nocent} remains in the outer well, and a displacement off a restoring radius still crosses that radius. That path is the travel of the electron, and the square of the pattern remains the one fixed by the degree. In a covalent bond the places are fixed, and the electrons occupy the configuration. Among the ions of a metal the electrons have no such place (Section~\ref{sec:free}). Quantum mechanics already reads an orbital as a stationary pattern. The signed force carries that reading through to the places themselves.

The diameter's slide toward a single seat is the overlap of Section~\ref{sec:pauli}. The gravitational pole expels it. The locus on which the signed forces cancel, and the Hessian is positive, is the two-center square.

\section{The laboratory well}
\label{sec:lab}

The sign that scattering and spectroscopy measure at atomic distances is the sign of the outer well. There \(\gamma_s>0\): a nucleus attracts an electron, electrons repel one another, and a path from any finite radius out to infinity does not cross a node. The integral~\eqref{eq:potential} therefore assigns a finite removal energy. The reversed layer does the opposite. Like charges attract, unlike charges repel, and the path to infinity is cut by the pole.

The restoring radii of Section~\ref{sec:shells} lie in that reversed layer. They are not ionization equilibria. For the cube the obstruction is geometric and holds for every nuclear charge \(Z\ge 3\). Each of the three chords---the edge, the face diagonal, and the body diagonal---is longer than the radius, so each electron--electron phase is smaller than the nuclear phase. On \((0,x_1)\) the interference factor is strictly decreasing, and the response built from those three chords therefore lies strictly below its far-field value. That value is the sum of the outward coefficients,
\[
\frac{3\sqrt{3}}{4}+\frac{3\sqrt{6}}{8}+\frac14,
\]
which lies strictly between \(2\) and \(3\). For \(Z\ge 3\) the response never meets the charge, and the outward force remains inward throughout the outer well. A neutral cube, and every heavier one, does not close where the laboratory signs hold. The same is already true of the diameter for every \(Z\ge 1\): the single coefficient is \(1/4\). That quarter is the whole of the two-electron problem. In any placement, while every distance remains in the outer well, the two electrons are not simultaneously force-free once the nuclear charge is at least \(1/4\). On opposite rays the repulsion stays strictly below a quarter of the attraction of unit nuclear charge at every finite radius. The electrons still spread: where the interference factor of their separation is positive, the antipode is an angular minimum, and the radius is not restored. The shell of two in this well is the two monopole seats of Section~\ref{sec:shells}.

The tetrahedron, the octahedron, and the icosahedron meet the same obstruction. Every chord is longer than the radius, so on \((0,x_1)\) each ratio of interference factors lies strictly below one and the response lies strictly below the far-field coefficient. For the tetrahedron that coefficient is \(3\sqrt{6}/8\), below \(1\). For the octahedron it is \(\sqrt{2}+1/4\), below \(2\): four edges of scale \(1/\sqrt{2}\) and one opposite vertex of scale \(1/2\). For the icosahedron it is \(\frac54\sqrt{5+2\sqrt{5}}+\frac14\), below \(5\): from each vertex, five neighbors, five farther vertices, and one opposite vertex. A nuclear charge at least \(1\), \(2\), or \(5\) therefore never meets the response, and the outward force remains inward through the outer well. A neutral shell of the tetrahedron, the octahedron, the cube, or the icosahedron does not close where the laboratory signs hold. The tetrahedron and the octahedron, like the diameter and the cube, each close once in the reversed layer, and the outward force restores the radius.

A restoring radius appears in the outer well when some like-charge chord is shorter than the radius. Shrinking the figure then drives that chord into the first node while the radius is still outside it, and the pair repulsion diverges outward. At large separation the same figure, if its geometric coefficient lies below the nuclear charge, falls inward. Between those two regimes the response crosses the charge.

Among the five regular polyhedra, only the dodecahedron has such a chord. From each vertex, three neighbors lie at squared distance \(2(3-\sqrt{5})/3\) times the squared radius, so the scale \(R/d\) of that chord is greater than one. The far-field coefficient of the twenty vertices lies strictly between \(15/2\) and \(8\). Before the nearest chord meets the first node the response stays strictly above \(7\). A nuclear charge of at most \(7\) therefore never meets the response while every pair is still in the outer well, and the outward force points outward on that whole interval. For every nuclear charge \(Z\ge 8\), the coefficient lies below the charge, so at large separation the force points inward. The response---the sum of the chord coefficients times \(\gamma_s\) at the radius divided by \(\gamma_s\) at the chord---equals \(Z\) at exactly one phase before the nearest chord meets the first node, and the radius there lies strictly between \(7\ell/10\) and \(10\ell/7\). A larger charge places that balance at a smaller radius. At a larger radius the force points inward, and at a smaller radius, while every pair is still in the outer well, it points outward. At the root the radial force vanishes, and the removal to infinity is the integral~\eqref{eq:potential}.

Twenty is the number of vertices, not a principal shell. The counts \(2\), \(8\), \(18\), and \(32\) remain a dual spinor and its spin-reversed partner on each mode. A shell of eight, packed as a cube, still closes only in the reversed layer. The dodecahedral root is a balance of twenty vertices in the outer well. It is not one of the four capacities.

\section{Electronegativity}
\label{sec:eneg}

In the outer well the attraction of a nucleus, counted positive and in units of \(k e^{2}/\ell^{2}\), is its charge times \(Y=4x^{2}/\gamma_s(x)\). At one common phase the factor \(Y\) is the same, and it is positive, so the larger charge pulls more strongly. \(Y\) is strictly increasing in the phase. A smaller charge therefore outpulls a larger one only by standing at a larger phase, which is a smaller radius.

Place charges \(Z\) and \(W\) at the ends of a segment of length \(D\), with an electron on the segment. Write \(\varphi=\ell/(2D)\) for the phase of the whole segment, and write \(t\) for the fraction of the segment lying between the electron and the charge \(Z\). Both distances lie in the outer well precisely when \(D\) is greater than twice the first-node radius and \(t\) lies strictly between \(\varphi/x_1\) and \(1-\varphi/x_1\). The net pull toward \(Z\) is \(Z\,Y(\varphi/t)-W\,Y(\varphi/(1-t))\). On that interval the pull is strictly decreasing in \(t\), and it vanishes at exactly one fraction.

At the midpoint the pull is \((Z-W)\,Y(2\varphi)\). When \(Z>W>0\) the electron there is drawn toward \(Z\), and the cancellation lies strictly closer to \(W\). On the side toward \(Z\) the force points toward \(Z\); on the side toward \(W\) it points toward \(W\). The cancellation is not restored, and it is not a seat. The part of the segment drawn toward the larger charge is the longer part. When the charges are equal the cancellation is the midpoint, and the electron is drawn toward the nearer nucleus. Raising the charge at one end moves the cancellation toward the other end.

The direction of this pull is the electronegativity of the larger charge. No laboratory scale is read off. Screening is not introduced, and \(\ell\) is not derived.

\section{Exclusion as gravitational repulsion}
\label{sec:pauli}

The gravitational channel uses the same interference factor, sourced by the kinematic pair tied to inertial mass rather than by the pair tied to charge~\cite{dst-main}. For two electron masses the force that increases their separation is
\begin{equation}
\label{eq:grav}
F_G=-\frac{G m_e^{2}}{\gamma_{s,G}(r)\,r^{2}},
\end{equation}
with \(\gamma_{s,G}\) the equal-scale factor of the mass channel at \(\lambda_\pi\), the length of the gravitational force of every mass. Where \(\gamma_{s,G}>0\) the masses attract. Where \(\gamma_{s,G}<0\) they repel. The first zero of \(\gamma_{s,G}\) is simple, by the same derivative~\eqref{eq:gamma}. Approached from inside, \(F_G\) therefore has a simple pole and points outward.

The prefactor is small. The CODATA values give
\begin{equation}
\label{eq:ratioG}
\frac{G m_e^{2}}{k e^{2}}=2.401\times 10^{-43}.
\end{equation}
A simple pole is not removed by a small prefactor. Write the inward fractional displacement from the first gravitational node as \(r=r_1(1-\varepsilon)\) with \(\varepsilon\downarrow 0\). Then \(x-x_1=x_1\varepsilon+O(\varepsilon^{2})\) and
\[
\bigl|\gamma_{s,G}\bigr|
  =2x_1\cosh x_1\sin x_1\,\varepsilon+O(\varepsilon^{2}).
\]
The coefficient lies strictly between \(\frac{\pi}{2\sqrt{2}}\cosh\frac{\pi}{4}\) and \(2\cosh 1\sin 1\). Numerically it is \(2.2260\).
The outward force exceeds the scale \(k e^{2}/r^{2}\) of a nonsingular electromagnetic force at the same radius when \(\varepsilon<1.078\times 10^{-43}\). The electron mass fixes the width of that neighborhood. It does not decide whether the neighborhood exists.

This comparison assumes that the electromagnetic factor is not itself zero at the same radius. The two channels are distinct phase pairs, so a gravitational node is not required to be an electromagnetic node, and the square of the next section does not choose \(\lambda_\pi\). If the two factors vanished together, their ratio would remain~\eqref{eq:ratioG} and the electromagnetic piece would dominate. The exclusion read here is the gravitational pole at a node the charge channel does not share.

Placing both electrons on one orbital seat drives their separation through that node. The force~\eqref{eq:grav} is then outward and unbounded, so the shared seat is left at once. The main text proposes that this repulsive layer is the Pauli principle: fermions entering it cannot occupy one seat, without a prior appeal to quantum statistics. The seat they cannot share is a repeated ray of the dual spinor. The second seat on the same mode is the orthogonal ray of the spin-reversed partner, and the two together are the singlet. Both seats stand on that one pattern and carry the same square, the response of the pattern when the direction is turned. That proposal is adopted here for electrons. It is not a derivation of the spin-statistics theorem, and the length of the mass channel remains \(\lambda_\pi\).

\section{The square and the covalent force}
\label{sec:covalent}

Place protons of charge \(+e\) at \(z=\pm R/2\) and electrons of charge \(-e\) at cylindrical radius \(\rho\), opposite each other in the midplane. The three separations are the proton--proton distance \(R\), the electron--electron distance \(r_{ee}=2\rho\), and the four equal electron--proton distances
\[
r_{ep}=\sqrt{(R/2)^{2}+\rho^{2}}.
\]
No angular momentum is assigned. Between the nuclei the electrons occupy fixed places in one configuration, and equilibrium means that the axial force on each proton and the cylindrical force on each electron both vanish.

The axial force on the proton at \(+R/2\), counted positive when it increases \(R\), and the cylindrical force on one electron, counted positive when it increases \(\rho\), are
\begin{align}
F_z
  &=k e^{2}\left(
    \frac{1}{\gamma_s(R)\,R^{2}}
    -\frac{R}{\gamma_s(r_{ep})\,r_{ep}^{3}}
  \right),
\label{eq:Faxial}\\
F_{\rho}
  &=k e^{2}\left(
    \frac{1}{\gamma_s(r_{ee})\,r_{ee}^{2}}
    -\frac{2\rho}{\gamma_s(r_{ep})\,r_{ep}^{3}}
  \right).
\label{eq:Frho}
\end{align}
The first term of \(F_{\rho}\) is the electron--electron force. In the outer well it is a repulsion and points outward: that outward piece is electrostatic, and it occupies the place a centrifugal term would have occupied. In the repulsive shell \(\gamma_s(r_{ee})<0\), so the same term points inward, because like charges attract there.

The configuration in which the four charges are the vertices of a square, with the protons on one diagonal and the electrons on the other, is
\begin{equation}
\label{eq:square}
\rho=\frac{R}{2},
\qquad
r_{ee}=R,
\qquad
r_{ep}=\frac{R}{\sqrt{2}}.
\end{equation}
On this locus \(F_z\) and \(F_{\rho}\) reduce to one expression. With \(x=\ell/(2R)\),
\begin{equation}
\label{eq:squareF}
F_z=F_{\rho}
  =\frac{k e^{2}}{R^{2}}
  \left(
    \frac{1}{\gamma_s(x)}
    -\frac{2\sqrt{2}}{\gamma_s(\sqrt{2}\,x)}
  \right).
\end{equation}
Both vanish when
\begin{equation}
\label{eq:squarecond}
\gamma_s(\sqrt{2}\,x)=2\sqrt{2}\,\gamma_s(x).
\end{equation}
The coefficient \(2\sqrt{2}\) is greater than one, and on \([0,\pi]\) the factor \(\gamma_s\) decreases. While the diagonal phase \(x\) is still less than the first node, \(\gamma_s(x)\) is positive and \(\sqrt{2}\,x\) has not yet reached \(\pi\). The value at the larger phase is then smaller than the value at \(x\), and that value is smaller than \(2\sqrt{2}\) times itself. The two sides of the balance do not meet. They still do not meet if only the side has crossed the node, for then one factor is negative and the other is positive. A square whose diagonal lies in the outer well is not an equilibrium.

The ratio of the electron separation to the proton separation may be written \(s\), so that the square is the case \(s=1\). While both separations lie in the outer well the two forces do not vanish together. Vanishing together would require the interference factor at the proton separation, divided by the factor at the electron separation, to equal \(s^{3}\). In that well the factor is positive and strictly decreasing in the phase. For \(s>1\) the electron separation has the smaller phase, so that ratio is strictly less than one, while \(s^{3}\) is strictly greater. For \(s<1\) the two comparisons exchange their directions. For \(s=1\) the ratio equals one, and the square condition stays strictly negative. Opposite signs of the two factors likewise fail to equal the positive value \(s^{3}\). With all three factors nonzero, the two forces vanish together only when \(\gamma_s(R)\,R^{3}\), \(\gamma_s(r_{ee})\,r_{ee}^{3}\), and \(\gamma_s(r_{ep})\,r_{ep}^{3}\) are one value. While the electron--proton separation lies in the outer well, no such value exists. On that well the quotient of the factor by the cube of the phase \(\ell/(2r)\) falls as the phase rises, so two phases there share the quotient only when they coincide, and the midplane then places the remaining phase in the same well at a different value. If neither of the other two phases lies in the well, each is at least the first node, and the electron--proton phase is at least \(\sqrt{2}\) times that node. A laboratory bond is a configuration in which the proton attracts the electron, so that separation lies in the outer well, and the pair has no equilibrium while it does.

The balance can hold only after the diagonal has entered the first repulsive shell, so that both the diagonal and the side lie between the two nodes. There both factors are negative: like charges attract along the diagonals, and unlike charges repel along the sides. On that range
\[
\gamma_s(\sqrt{2}\,x)-2\sqrt{2}\,\gamma_s(x)
\]
increases with \(x\). It is negative at \(3/2\) and positive at \(8/5\), so it vanishes once, and only once, in the shell. The crossing lies between those two phases. At \(8/5\) the side phase is still less than \(\pi\), and the second node lies beyond \(\pi\), so the crossing is inside the shell. On the bracket between \(3/2\) and \(8/5\) both factors stay negative, and the common force on the square has the sign of the balance. A larger diagonal is drawn inward and a smaller one is pushed outward, so the root restores the square. Its phase is
\[
x=1.55463,
\qquad
\frac{R}{\ell}=0.32162,
\qquad
\frac{r_{ep}}{\ell}=0.22742.
\]
The two interference factors are \(\gamma_s(x)=-2.2208\) on the diagonals and \(\gamma_s(\sqrt{2}\,x)=-6.2812\) on the sides. The repulsion that holds the electrons off the protons is the side force. It is not a centrifugal force.

The partial forces at this root are
\begin{align}
\bigl|F_{\mathrm{diag}}\bigr|
  &=\frac{k e^{2}}{\bigl|\gamma_s(x)\bigr|\,R^{2}}
  =4.3532\,\frac{k e^{2}}{\ell^{2}},
\label{eq:Fdiag}\\
\bigl|F_{\mathrm{side}}\bigr|
  &=\frac{k e^{2}}{\bigl|\gamma_s(\sqrt{2}\,x)\bigr|\,r_{ep}^{2}}
  =3.0782\,\frac{k e^{2}}{\ell^{2}}.
\label{eq:Fside}
\end{align}
The diagonal force is the covalent force scale: it is the attractive force along the proton--proton separation, cancelled by the axial part of the side repulsions. The side force is the diagonal force divided by \(\sqrt{2}\).

The electrostatic potential, regarded as a function of \(R\) and \(\rho\) near this root, has Hessian eigenvalues \(37.33\,k e^{2}/\ell^{3}\) and \(7.933\,k e^{2}/\ell^{3}\), both positive. The root is a local minimum. Allowing \(\rho\) to follow \(R\), the stiffness of the diagonal is
\begin{equation}
\label{eq:kappa}
\kappa=8.178\,\frac{k e^{2}}{\ell^{3}}.
\end{equation}
Two further roots of \(F_z=F_{\rho}=0\) lie in the same shell. Each has Hessian eigenvalues of opposite sign, so each is a saddle. The square is the stable point of the family. Its two electrons occupy opposite vertices, a separation \(R\), not one seat. A displacement that puts both on one vertex is the overlap of Section~\ref{sec:pauli}: once that separation reaches the first gravitational node, \(F_G\) expels them. The electromagnetic Hessian stabilizes the square against small displacements. The gravitational pole forbids the shared seat even though \(m_e\) is small. Because the root lies inside the shell, the outer-well integral~\eqref{eq:potential} does not assign it a dissociation energy relative to separated atoms at infinity.

Helium is a nucleus of charge two with two electrons, the one-center shell of two. Place charge two on each nucleus of the same midplane family and keep the two electrons. On the square the axial balance asks for the interference factor on the side to be \(\sqrt{2}\) times the factor on the diagonal, and the cylindrical balance asks for \(4\sqrt{2}\) times that factor, so the square is not an equilibrium wherever both factors are finite and nonzero. While the internuclear separation and the electron separation both lie in the outer well, the two forces do not vanish together at any ratio. Where the electrons are farther apart than the nuclei, the ratio of the factors is strictly less than one, and the balance asks for four times the cube of the ratio of the separations, with that cube strictly greater than one. Where the electrons are closer, each electron lies closer to a nucleus than the nuclei lie to each other. The axial balance then asks for more than \(\sqrt{2}\) times the factor on the internuclear separation, but the phase of the shorter separation stays below the second node, so the factor there is smaller when it is positive and negative when the phase has passed the first node. Opposite signs of the factors at the two separations fail as well, because the balance asks for a positive ratio. The laboratory sign is the nucleus attracting the electron, and this pair has no equilibrium in that well. Charge two with two electrons remains the one-center shell of two, the two monopole seats.

\section{The water molecule}
\label{sec:water}

The water molecule places the nucleus at the origin and one proton on each of two rays, with one bonding electron on each ray. The electron lies closer to the nucleus than its proton. Further electrons may stand off the rays. The angle between the rays may vary, and the electron and the proton on one ray stay distinct.

While the two bonding electrons repel, the other bond drives each bonding electron toward a wider angle. The proton on the other ray pulls the other way, and that pull does not win. The nuclear charge does not enter this comparison. An electron on the same ray adds nothing to the tangential force. An electron that repels the bonding electron from the side of the ray toward the other bond drives it toward a still wider angle, and one that repels it from the opposite side drives it toward a narrower angle. So long as every further electron stands on the ray or repels the bonding electron from the side toward the other bond, the tangential force on the bonding electron stays toward the wider angle. While the bonding electrons repel and lie closer than their protons, that force is positive at every angle short of a straight line and vanishes on the line. The straight line restores the bonding electrons.

The tangential force on each proton is another balance. The nucleus and the electron on the same ray contribute nothing to it, and the nuclear charge does not enter. What remains is the other proton's repulsion minus the pull of the other bonding electron. While both separations lie in the outer well and the protons are no farther apart than each is from the other electron, the force points toward the wider angle. The angle between the rays at which the two separations are equal is strictly less than a third of a straight angle. When the electron lies at least a quarter as far from the nucleus as the proton, and that equal separation still lies in the outer well, the same difference is negative on the straight line. The force on the proton is then positive at the equal separation and negative at some greater angle still short of the straight line, so it vanishes at an angle strictly between them. The straight line does not restore the protons. If the separation of the bonding electrons is still in the outer well at that angle, their tangent points toward the wider angle, and they are not at rest on the ray where the force on the proton vanishes.

On a straight line, with both distances from the nucleus in the outer well, the two bonds alone exert no tangential force. An electron on the axis opposite the bonds then takes the arrangement off the line. The bonds still drive each bonding electron toward a wider angle at every shorter angle where those electrons repel. While the opposite electron's separation lies in the outer well, that electron drives them toward a narrower angle, and on the straight line the tangent is strictly negative, so the straight line no longer restores them.

On that straight line, place each electron between the nucleus and its proton, with both distances from the nucleus in the outer well and with the interference factor on the segment between them not a node. The outward force on the electron and the outward force on the proton then sum to a strictly negative value whenever the nuclear charge is at least \(1/4\). The attraction along that segment cancels in the sum. The molecule is these two bonds together with the electrons that stand off them.

\section{The angle of a covalent bond}
\label{sec:angle}

Each ray from the nucleus carries the electron of one pair. The proton of a shared pair lies farther out on that ray and draws the electron outward. A non-shared pair has no proton on its ray. Place four electrons at one distance on the rays of a regular tetrahedron, and protons at one greater common distance on some of those rays, with every separation in the outer well. The outward force is then strictly greater on an electron whose ray carries a proton than on an electron whose ray does not. The non-shared electron is driven inward and is the closer one.

The angle between two of those rays has cosine \(-1/3\) and is greater than a right angle. Four electrons at one distance and four protons at one greater distance exert no tangential force. That balanced angle is the angle of four shared pairs. While an angle is at least a right angle and both separations lie in the outer well, moving an electron closer shortens its separation from a neighbor, the phase rises, and the repulsion \(Y=4x^{2}/\gamma_s\) strengthens. If the closer electron is still at least \(1/\sqrt{2}\) times as far from the nucleus as the neighbor, the tangential part of that repulsion is stronger as well.

Keep three of the electrons at the common distance. Move the fourth inward, no closer than \(1/\sqrt{2}\) times the neighbor's distance from the nucleus, and remove the proton on its ray, with the separations still in the outer well. At cosine \(-1/3\) the tangential force on each of the three points away from the electron that moved in, and the angle between those three is driven narrower. A second electron moved inward by the same amount, with its proton likewise removed, contributes the same narrowing, so the two contribute twice the narrowing of one. A non-shared pair that sits closer to the nucleus therefore pushes the bonds aside more strongly than a shared pair at the bonding distance. Two such pairs push twice as hard as one, and the force on the remaining bonds drives their angle below the angle of four equal pairs. On two rays alone, with nothing off them, the straight line of Section~\ref{sec:water} still restores the bonding electrons. The two bonds of water are two shared pairs with two closer electrons standing off those rays.

\section{Coordination and the metallic force}
\label{sec:metal}

A nearest-neighbor bond in a monovalent coordination shell of \(z\) neighbors is given the electron content \(\nu=2/z\). The dimer is \(z=1\), \(\nu=2\), the square of Section~\ref{sec:covalent}. The charge \(-\nu e\) is split into two lumps of charge \(-(\nu/2)e\) at opposite points of the midplane. The forces~\eqref{eq:Faxial} and~\eqref{eq:Frho} become
\begin{align}
F_z
  &=k e^{2}\left(
    \frac{1}{\gamma_s(R)\,R^{2}}
    -\frac{\nu R}{2\,\gamma_s(r_{ep})\,r_{ep}^{3}}
  \right),
\label{eq:Fznu}\\
F_{\rho}
  &=k e^{2}\left(
    \frac{(\nu/2)^{2}}{\gamma_s(r_{ee})\,r_{ee}^{2}}
    -\frac{\nu\rho}{\gamma_s(r_{ep})\,r_{ep}^{3}}
  \right).
\end{align}
Equilibrium is still \(F_z=F_{\rho}=0\), with no centrifugal term. The outward piece in the outer well is the first term of \(F_{\rho}\): the lumps repel where \(\gamma_s(r_{ee})>0\).

For \(z=8\) and \(z=12\) that outer-well system has one root each. All three separations lie outside \(r_1\), so the signs are the ordinary ones: like charges repel, unlike charges attract. The partial proton--proton force and its ratio to the covalent diagonal~\eqref{eq:Fdiag} are
\begin{align}
z=8:\quad
\frac{R}{\ell}&=2.3591,&
F_{\parallel}&=0.18820\,\frac{k e^{2}}{\ell^{2}},&
\frac{F_{\parallel}}{|F_{\mathrm{diag}}|}&=0.04323,
\label{eq:ratio8}\\
z=12:\quad
\frac{R}{\ell}&=1.2666,&
F_{\parallel}&=0.74188\,\frac{k e^{2}}{\ell^{2}},&
\frac{F_{\parallel}}{|F_{\mathrm{diag}}|}&=0.1704.
\label{eq:ratio12}
\end{align}
At both roots the Hessian of the potential has eigenvalues of opposite sign. Each is a saddle, not a local minimum. The same is true of the further roots that sit inside the repulsive shell. Reducing the electron content from a full pair to \(2/z\) leaves a point where the signed forces cancel, and it does not leave a stable configuration. The saddles are not a defect of the lumps. No point in the outer well of every ion holds an electron (Section~\ref{sec:free}). An empirical nearest-neighbor distance would fix \(F_{\parallel}\) by the proton--proton term alone; it is an input, and it is not required for the ratios.

\section{Free electrons}
\label{sec:free}

The electrons of a metal stand where every ion attracts them, in the outer well of each ion. Displace such an electron from a point along a direction. The component of the force along that direction changes at a definite rate, and a positive rate means that the displacement is pushed further. For one point charge \(Q\), counted in units of \(e\), at phase \(x=\ell/(2r)\), the rates along any three perpendicular directions sum to
\begin{equation}
\label{eq:outflow}
Q\,\frac{k e^{2}}{\ell^{3}}\,S(x),
\qquad
S(x)=\frac{16x^{4}\cosh x\sin x}{\gamma_s(x)^{2}}.
\end{equation}
On the outer well \(S\) is positive, and it strictly increases with the phase. An ion therefore contributes a positive sum, and a further electron a negative one. With the Coulomb force the same sum is zero. That zero is Earnshaw's theorem: under the Coulomb law no arrangement of charges holds a further charge in every direction. The signed factor removes the zero, and it removes it in opposite directions for the two signs of the charge.

An electron in the outer well of every ion therefore has, in any three perpendicular directions, at least one along which the ions push a displacement further. Further electrons act the other way, and \(S\) grows as the distance falls. So long as each further electron lies in the outer well and farther away than a distinct ion, and each ion carries a charge of at least one, the ions still win. No point between the ions holds the electron. Where the forces on it cancel, the cancellation is not restored in every direction. Section~\ref{sec:eneg} found the same along a segment.

Symmetry makes the rate the same in every direction. At the center of a regular octahedron, a cube, and a cuboctahedron of ions, of which one or two may be absent, the force on an electron vanishes, and every displacement is pushed further at one rate. In units of \(k e^{2}/\ell^{3}\), each shell contributes a third of its number of ions, \(6\), \(8\), or \(12\), times its charge times \(S\) at its phase. The nearest two shells about a site of the simple cubic, the body-centered cubic, and the face-centered cubic arrangement are of these kinds. So are the nearest two shells about the octahedral hole of the face-centered cubic arrangement, an octahedron and then a cube.

The same shells hold a positive ion at their center. The force on the ion vanishes, and every displacement of the ion is restored, at the rate of the electron times the charge of the ion, with the sign reversed. Under the Coulomb law both rates are zero. In the outer well the ions are held in place by their like neighbors, and the electrons are pushed off every such place. A metal in the laboratory well is a held arrangement of ions among which the electrons have no seat. The electrons are free. The covalent square of Section~\ref{sec:covalent} is the opposite case: its places lie in the reversed shell, and its Hessian is positive definite.

No infinite lattice sum is taken, and no conductivity or cohesive energy is derived. Screening is not introduced: each ion is a point charge. Section~\ref{sec:crystal} counts every shell of a cubic truncation.

\section{Crystals}
\label{sec:crystal}

The rate \(S\) of~\eqref{eq:outflow} has the sign of \(\sin x\). It is positive at every phase in \((0,\pi)\) off the first node, on both sides of that node. It vanishes at \(\pi\), and it is negative on \((\pi,2\pi)\) wherever \(\gamma_s\) is not zero. Like charges about a symmetric site therefore hold one another on both sides of the node, out to the radius \(\ell/(2\pi)\), and push one another once the phase exceeds \(\pi\). The reversal is that radius, not the node. With \(\gamma_s=1\) the three-direction sum is zero.

A finite set of point charges about a site, carried into itself, charges included, by the rotations of a regular tetrahedron, exerts no force at the site. The three-fold turn about a body diagonal and the half-turn about a coordinate axis are enough. A displacement of a test charge \(q\) then changes the outward component at one rate in every direction, \(-q/3\) times \(\sum Q\,S\). Any piece of the simple cubic lattice, the body-centered cubic lattice, or the face-centered cubic lattice, cut by a cube about the site and by any radius, is such a set. While every charge is positive and every phase lies below \(\pi\) and off the first node, the site holds a positive ion and pushes an electron off. The phase falls as the distance grows, so a nearest neighbor in the outer well places every farther shell there as well, and the holding is the whole piece, to any number of shells. No infinite sum is taken.

On the outer well, \(S(x)/x^{4}\) strictly increases. Consider a site whose nearest shell is unlike, \(N_{1}\) ions of charge \(Q_{1}\) at squared distance \(k_{1}\), and whose next shell is like, \(N_{2}\) ions of charge \(Q_{2}\) at squared distance \(k_{2}\), with the nearer shell in the outer well. The tetrahedral symmetry still cancels the force, and a positive test charge is pushed off in every direction once
\begin{equation}
\label{eq:twoshell}
N_{2}Q_{2}k_{1}^{2}\le N_{1}Q_{1}k_{2}^{2}.
\end{equation}
The nearest two shells of four arrangements meet the bound, and so do their ordinary charges. Rock salt has six unlike neighbors at squared distance \(1\) and twelve like neighbors at squared distance \(2\), and~\eqref{eq:twoshell} is \(Q_{2}\le 2Q_{1}\). Cesium chloride has eight at squared distance \(3\) and six at squared distance \(4\), and zincblende has four at squared distance \(3\) and twelve at squared distance \(8\); both bounds are \(27Q_{2}\le 64Q_{1}\). Fluorite, seen from the cation, has eight at squared distance \(3\) and twelve at squared distance \(8\), and the bound is \(27Q_{2}\le 128Q_{1}\). Seen from the anion it has four at squared distance \(3\) and six at squared distance \(4\), and the bound is \(27Q_{2}\le 32Q_{1}\). Equal charges meet the first three bounds. A cation of twice the anion's charge meets both sites of fluorite. These are the nearest two shells only. The Coulomb rate on the same sites is zero. The two shells do not hold the ion.

A linear map carrying a lattice into itself has an integer trace. A rotation through an angle \(\theta\), proper or improper, has trace \(2\cos\theta+\varepsilon\) with \(\varepsilon=\pm 1\). A rotation through \(2\pi/n\) therefore preserves a lattice in three dimensions only for \(n\in\{1,2,3,4,6\}\). A crystal site has no fivefold axis. The icosahedron and the dodecahedron of Section~\ref{sec:lab} are not the symmetry of a lattice site. As one-center shells they remain what that section computed. No cohesive energy and no conductivity are derived, and each ion is a point charge.

\section{A numerical illustration}
\label{sec:numbers}

Lengths in this section are multiples of \(\ell\), and forces are multiples of \(k e^{2}/\ell^{2}\). The electromagnetic length is the reduced Compton wavelength of the electron.

The laboratory sign at a one-electron radius is an input, and it places this length under a cap. Attraction at the Bohr radius \(a_0\) places that radius in the outer well, so \(\ell<2x_1 a_0\). The first node lies strictly below \(19/20\), and the cap is therefore \(\ell<(19/10)a_0\). On the SI stand-ins that quantity lies strictly between \(1005\times 10^{-13}\,\mathrm{m}\) and \(1006\times 10^{-13}\,\mathrm{m}\). The same sign at the hydrogenic radius \(a_0/Z\), while the charge \(Z\) still attracts the electron, tightens the cap to \(\ell<(19/10)a_0/Z\). A larger charge makes the cap smaller. For the input \(Z=2\) the cap lies strictly between \(502\times 10^{-13}\,\mathrm{m}\) and \(503\times 10^{-13}\,\mathrm{m}\). For the input \(Z=92\) it lies strictly between \(1092\times 10^{-15}\,\mathrm{m}\) and \(1093\times 10^{-15}\,\mathrm{m}\). Placing the first node at \(a_0\) makes the force diverge at that radius, and the phase of the radius \(a_0/2\) then lies where \(\gamma_s<0\). Attraction at \(a_0/2\) leaves that placement outside the outer well.

The electromagnetic length is the reduced Compton wavelength of the electron, \(\hbar/(m_e c)\). On the SI stand-ins it lies strictly between \(3861\times 10^{-16}\,\mathrm{m}\) and \(3862\times 10^{-16}\,\mathrm{m}\), and strictly below the charge-\(92\) cap. Its first node lies strictly between \(10/19\) and \(15/28\) of that length, hence strictly between \(2032\times 10^{-16}\,\mathrm{m}\) and \(2069\times 10^{-16}\,\mathrm{m}\), and that node lies strictly inside \(a_0/92\). The radii \(a_0\), \(a_0/2\), and \(a_0/92\) stay in the outer well. The caps lie above this wavelength.

The circular radius of Section~\ref{sec:nocent}, the first node, the diagonal of the restoring square, its side, the one-electron extremum, and the second node stand in that order, each outside the next. The circular comparison changes at a phase strictly between \(1/2\) and \(3/5\), so that radius lies strictly between \(5\ell/6\) and \(\ell\). The first node lies strictly between \(14/15\) and \(19/20\), and its radius strictly between \(10\ell/19\) and \(15\ell/28\), with \(15\ell/28\) strictly less than \(5\ell/6\). The square of Section~\ref{sec:covalent} lies strictly between \(153/100\) and \(157/100\), so the diagonal lies strictly between \(50\ell/157\) and \(50\ell/153\), with \(50\ell/153\) strictly less than \(10\ell/19\). The side is the diagonal divided by \(\sqrt{2}\), and that length lies strictly below \(50\ell/157\) and strictly above \(\ell/5\). The extremum lies strictly between \(\ell/(2\pi)\) and \(\ell/5\). The second node lies in \((\pi+\pi/4,\pi+\pi/2)\), and its radius lies strictly between \(\ell/(3\pi)\) and \(2\ell/(5\pi)\), with \(2\ell/(5\pi)\) strictly less than \(\ell/(2\pi)\). The values \(x_1=0.93755\), \(r_1/\ell=0.53330\), the phase \(1.55463\) with diagonal \(0.32162\,\ell\), and \(r_2/\ell=0.12731\) lie in these intervals.

The diagonal force at the square is \(4x^{2}/\bigl|\gamma_s(x)\bigr|\), and on \((153/100,157/100)\) it is strictly greater than \(4\). At \(x=\pi\), where \(\gamma_s=-\cosh\pi\), the outward force on an unlike pair is \(4\pi^{2}/\cosh\pi\), strictly greater than \(3\) and strictly less than \(7/2\). The covalent diagonal therefore exceeds the one-electron force at the extremum of \(\gamma_s\). The side force is the diagonal force divided by \(\sqrt{2}\), and is strictly greater than \(2\sqrt{2}\). The values \(4.3532\), \(3.4057\), and \(3.0782\) meet these bounds.

At atomic distances the measured sign is attraction, so the Bohr radius lies in the outer well, where a proton attracts an electron. The square lies in the reversed shell, and along the side the proton repels the electron. Placing that side at the Bohr radius would put the first node at \(\sqrt{2}\,x/x_1\) Bohr radii, strictly between \(91/40\) and \(12/5\). The laboratory sign forbids the identification. No dissociation energy is quoted: the equilibrium lies in the repulsive shell, and~\eqref{eq:potential} does not cross the node.

\section{Vibration across the sign boundary}
\label{sec:heat}

A restoring radius is the boundary between an inward force and an outward force. On the side of larger radius the outward force points inward, and on the side of smaller radius it points outward. A displacement off the root oscillates across that change of sign. The node is not the seat of the oscillation. At the node the force is undefined, and the path below stays where \(u\) is finite.

The outward force on one electron of a symmetric shell is
\begin{equation}
\label{eq:Fout}
F_{\mathrm{out}}(r)=\frac{k e^{2}}{r^{2}}\,u(x),
\qquad
x=\frac{\ell}{2r},
\end{equation}
with \(u\) the combination of nuclear and chord terms already used for the diameter and for the dodecahedron. At a simple root \(u(x_*)=0\) with \(u'(x_*)>0\),
\begin{equation}
\label{eq:curv}
F_{\mathrm{out}}'(r_*)=-\frac{k e^{2}}{\ell^{3}}\,K,
\qquad
K=8x_*^4\,u'(x_*).
\end{equation}
The electron mass turns the curvature into a frequency,
\begin{equation}
\label{eq:omega}
\omega^{2}=\frac{K\,k e^{2}}{m_e\ell^{3}}.
\end{equation}
Inside this linear reading \(\omega\) is fixed by \(K\), by \(m_e\), and by \(\ell\). It does not grow with the amplitude, and it does not grow with the temperature.

The radial root leaves the angular velocity free. In polar coordinates
\begin{equation}
\label{eq:polar}
m_e\bigl(\ddot r-r\dot\theta^{2}\bigr)=F_{\mathrm{out}}(r).
\end{equation}
The term \(m_e r\dot\theta^{2}\) is the kinematics of a tangential velocity. It is not a new contribution in~\eqref{eq:force}. While the velocity is small the root shifts outward by
\begin{equation}
\label{eq:shift}
\delta r=\frac{m_e\langle v_{\theta}^{2}\rangle}{|F_{\mathrm{out}}'(r_*)|\,r_*}.
\end{equation}
When the tangential energy equals the energy of a radial amplitude \(A\), the shift is \(A^{2}/r_*\), smaller than \(A\) by the factor \(A/r_*\). The mean radius remains the electrostatic root. The electron goes around the shell while the radius oscillates across the sign boundary. For the diameter the angle is a maximum, by Section~\ref{sec:shells}, so that going-around is not itself a vibration. The vibration is the radial crossing.

Temperature is the energy in the oscillator. Classical equipartition of one radial mode gives
\begin{equation}
\label{eq:amp}
\langle(\delta r)^{2}\rangle
  =\frac{k_B T}{|F_{\mathrm{out}}'(r_*)|}
  =\frac{k_B T\,\ell^{3}}{K\,k e^{2}}.
\end{equation}
The root-mean-square amplitude grows as \(\sqrt{T}\). That growth is the vigor of the motion. The constants \(k_B\) and, below, \(c\) and \(h\) are laboratory constants. They are not fixed by \(\gamma_s\).

An accelerated charge radiates. The mean Larmor power of one nonrelativistic electron is
\begin{equation}
\label{eq:larmor}
P=\frac{2k e^{2}\langle a^{2}\rangle}{3c^{3}},
\qquad
a=-\omega^{2}\,\delta r.
\end{equation}
With~\eqref{eq:amp}, \(\langle a^{2}\rangle=\omega^{2} k_B T/m_e\), so \(P\) is proportional to \(T\) at fixed \(\omega\), and proportional to \(K\) at fixed \(T\). A hotter occupation of the same root radiates more power at the same frequency. A warm body is a large collection of roots. Its heat rays are this radiation. They are not the line of one root.

The curvatures already in hand are far apart. The charge-\(8\) dodecahedron, the unique twenty-vertex root of Section~\ref{sec:lab}, has \(K=4.556\). The diameter of charge \(1\), in the reversed layer, has \(K=126.45\). On the square, at fixed protons, the curvature in \(\rho\) is \(18.10\,k e^{2}/\ell^{3}\), and~\eqref{eq:omega} applies with that number in place of \(K\). The diagonal stiffness~\eqref{eq:kappa} is carried by the proton reduced mass \(m_p/2\).

The classical power favors the stiffer root, because \(P\) grows with \(K\). The heat ray of a warm body is the enclosure spectrum read below, and a spacing large compared with \(k_B T\) leaves an oscillator empty. The laboratory displacement law puts the peak of a body at \(300\,\mathrm{K}\) at \(9.659\,\mu\mathrm{m}\), where \(k_B T=0.02585\,\mathrm{eV}\)~\cite{codata2018}.

The frequency in the Larmor formula is fixed by the sign boundary before a temperature is assigned. With the inertia of the moving coordinate in place of \(m_e\), \eqref{eq:omega} is the whole of it: the dimensionless curvature of the root, that inertia, and \(\ell\). Temperature is the energy with which the root is occupied. It determines the amplitude~\eqref{eq:amp}, and through \(a=-\omega^{2}\,\delta r\) the power~\eqref{eq:larmor}, at a frequency the occupation does not retune.

One signed balance can still radiate at two frequencies, because two coordinates can move on it. On the square the cylindrical motion of an electron at fixed protons has curvature \(18.10\,k e^{2}/\ell^{3}\) and inertia \(m_e\). The diagonal has stiffness \(8.178\,k e^{2}/\ell^{3}\) and inertia \(m_p/2\). By~\eqref{eq:omega} the ratio of the wavelengths is the square root of the ratio of \(K/m\) on the electron coordinate to \(K/m\) on the diagonal. That ratio is \(45.08\). Both wavelengths scale as \(\ell^{3/2}\), so the ratio does not use an identification of the side. The electron is accelerated at the nuclear period because the side is constrained to the diagonal, while its own crossing of the sign boundary keeps the curvature \(18.10\,k e^{2}/\ell^{3}\).

Temperature moves the occupation along the order those roots already form. The order is the order of \(K/m\). At \(300\,\mathrm{K}\) the displacement peak lies at \(0.1283\,\mathrm{eV}\), which is \(4.96\,k_B T\). Raising the temperature raises \(k_B T\), and a higher root of the signed force can enter the occupation. The vigor that belongs to ordinary temperature is the slow end of the order, the soft boundary. The lines themselves stay fixed. The displacement of the continuous peak is the occupation of a ladder whose spacing is measured in units of \(k_B T\).

A transition has no such motion inside its frequency. The condition
\begin{equation}
\label{eq:bohrnu}
\nu=\frac{E_i-E_j}{h}
\end{equation}
takes \(\nu\) from a pair of eigenvalues. The eigenvalues are the spectrum of a Hamiltonian. Temperature parameterizes the ensemble over that spectrum and leaves the spectrum unchanged, so the difference is the same in a cool body and in a hot one. What changes is the rate, through the occupation of the upper state. A difference fixed by \(K/m\) and by \(\ell\) is invariant in that sense. The peak of the heat ray is a different object. It sits at \(4.96\,k_B T\) and scales with the temperature. To identify the peak with a transition is to require one frequency to be both an eigenvalue difference and proportional to \(T\). The spectrum of a fixed Hamiltonian admits no such proportion. That is the contradiction.

The heat ray is the blackbody radiation of a warm body. A stationary orbital fixes the electron's motion. The proton carries the thermal velocity of the atom. The difference of the two velocities is not a quantity the orbital computes, and the difference does not change with the temperature. The wavelength of the blackbody changes with the temperature. The orbital picture therefore does not explain the blackbody.

The occupation of one such line is a geometric series. Write the spacing as \(x=\varepsilon/(k_B T)\) with \(x>0\). The mean number of quanta above the ground state is
\[
\langle n\rangle=\frac{1}{e^{x}-1},
\]
and the mean energy above the ground state is \(\varepsilon\langle n\rangle\). As \(x\) tends to zero this energy returns to \(k_B T\). Once \(x\ge 1\) the mean number of quanta lies below \(2e^{-x}\). Where the spacing is small, the classical power at fixed temperature still grows with \(K\), and a stiffer root radiates more. Where the spacing is large, the same weight empties the oscillator. The spacing remains the frequency~\eqref{eq:omega}.

A large enclosure sums many such lines, weighed by the number of modes that share a frequency. In a cavity the number of modes in an interval \(d\nu\) grows as \(\nu^{2}\,d\nu\). That growth is the geometry of the cavity. Planck's \(h\) remains the external constant that converts a frequency into the spacing \(x=h\nu/k_B T\), as \(k_B\) and \(c\) remain external. The energy in \(d\nu\) is then proportional to \(x^{3}/(e^{x}-1)\). This shape has a single positive maximum, at the unique positive solution of \((3-x)e^{x}=3\), and that solution lies between \(5/2\) and \(3\). The frequency of the maximum is proportional to the temperature.

Counted by wavelength, \(\nu=c/\lambda\) contributes the factor \(c/\lambda^{2}\). The shape becomes \(x^{5}/(e^{x}-1)\), with \(x=hc/(\lambda k_B T)\). Its single maximum lies between \(9/2\) and \(5\), at the unique positive solution of \((5-x)e^{x}=5\). The laboratory peak \(4.96\,k_B T\) is this wavelength maximum. The frequency maximum is the smaller multiple, below \(3\,k_B T\). The value of \(x\) that locates the maximum does not depend on \(T\), so the product of the peak wavelength and the temperature is constant. Temperature changes the unit in which one fixed ladder is read. That constant product is Wien's displacement.

The density proportional to \(T\nu^{2}\), the law of Rayleigh--Jeans, increases at every frequency, so equipartition on the cavity has no maximum. The maximum appears when the geometric weight \(1/(e^{x}-1)\) stands in place of the share \(k_B T\).

The integral of the frequency shape is
\[
\int_{0}^{\infty}\frac{x^{3}}{e^{x}-1}\,dx=\frac{\pi^{4}}{15}.
\]
The integrated frequency density of the enclosure therefore scales as \(T^{4}\). The radiated power of one classical root scales as \(T\). This fourth power is Stefan--Boltzmann's law for that integral. The factors that convert it into a flux, among them the two polarizations of the cavity, remain outside the count. The heat ray of the warm body is the integral. The restoring roots remain an order of fixed frequencies, and each keeps~\eqref{eq:omega}. Ordinary temperature still occupies the heavy end of that order. A collision may move an electron from one place on a sign boundary to another, and the field radiated there is the Larmor field of that root. The weight \(1/(e^{x}-1)\) decides how far up the order a contribution survives. The cavity carries the continuous spectrum.

\section{Conclusions}

The sign of \(\gamma_s\) is the repulsion in~\eqref{eq:force}. A centrifugal term is not required, and it is not introduced. One electron and one proton have no force-free radius where the two interference factors have the same sign. At the Bohr radius the mass term is smaller than the charge term by a factor strictly between \(4.40\times 10^{-40}\) and \(4.42\times 10^{-40}\). In the outer well a circular path is radially restoring above a radius that lies strictly between \(5\ell/6\) and \(\ell\), and unstable between that radius and the first node. The circular paths of that well have their least energy at the same radius, and the value is negative. One electron about a nucleus of any positive charge keeps that radius, and the energy of the path is the charge times the one-proton value. A diameter, a tetrahedron, an octahedron, and a cube about one nucleus each have one restoring radial balance in the first repulsive shell, the diameter is an angular maximum, and \(18\) and \(32\) count a dual spinor and its spin-reversed partner through the quadrupole and the octupole rather than a node. The square of the angular momentum on the monopole, the dipole, the quadrupole, and the octupole, their response when the direction is turned, is \(0\), \(2\), \(6\), and \(12\). A seat is a pattern on directions. What is called an electron orbital is an electron configuration, not an orbit. An electron may still travel, and in a covalent bond the places are fixed. Only the monopole is independent of direction, and every higher seat is positive in some directions and negative in others. Two electrons in the outer well are not simultaneously force-free once the nuclear charge is at least \(1/4\): on opposite rays the repulsion is strictly less than a quarter of the attraction of unit nuclear charge, and where the interference factor of the separation is positive the antipode is an angular minimum. The shell of two in that well is the two monopole seats, not a zero of the signed force. The tetrahedron, the octahedron, the cube, and the icosahedron have no root in the outer well once the nuclear charge is at least \(1\), \(2\), \(3\), and \(5\). A regular dodecahedron does. Its far-field coefficient lies strictly between \(15/2\) and \(8\), and before the nearest chord meets the first node the response stays strictly above \(7\), so a charge of at most \(7\) is pushed outward on that whole interval. For every \(Z\ge 8\) the response equals the charge at exactly one phase before that chord meets the node, the radius lies strictly between \(7\ell/10\) and \(10\ell/7\), a larger charge places that balance at a smaller radius, and the removal to infinity stays in the outer well. On a segment longer than twice the first-node radius, the electron is drawn from the midpoint toward the larger charge, and the one cancellation lies closer to the smaller charge and is not restored. Two of each have exactly one square in the first repulsive shell, and that square restores the diagonal. Its phase lies between \(3/2\) and \(8/5\), with diagonal force \(4.3532\,k e^{2}/\ell^{2}\) and stiffness \(8.178\,k e^{2}/\ell^{3}\). While the proton separation and the electron separation both lie in the outer well, the two forces do not vanish together at any ratio of those separations. While the electron--proton separation lies in the outer well, the two forces do not vanish together. Helium, a nucleus of charge two with two electrons, remains the shell of two: the same family with charge two on each nucleus has no equilibrium in that well, and the square is not an equilibrium wherever both interference factors are nonzero. The water molecule of Section~\ref{sec:water} keeps the tangential force on each bonding electron toward the wider angle while those electrons repel and lie closer than their protons, so long as every further electron stands on the ray or repels from the side toward the other bond, and the straight line restores them. On the straight line the outward sum is negative once the nuclear charge is at least \(1/4\). When the electron lies at least a quarter as far from the nucleus as the proton and the equal proton separation still lies in the outer well, the tangential force on each proton vanishes strictly between that separation and the straight line, and the angle between the rays at the equal separation is strictly less than a third of a straight angle. Four shared pairs on tetrahedral rays exert no tangential force at the angle of cosine \(-1/3\), which is greater than a right angle. A non-shared electron is driven inward because its ray has no proton, and while it stays at least \(1/\sqrt{2}\) times as far from the nucleus as a neighbor it pushes the remaining bonds aside more strongly than a shared pair at the bonding distance. Two such electrons contribute twice that narrowing. The bonds that carry \(2/8\) and \(2/12\) of an electron pair have electrostatic roots and indefinite Hessians; their partial forces are \(0.04323\) and \(0.1704\) of the covalent diagonal. In the outer well of every ion the ions alone hold no electron: the rates at which three perpendicular displacements are pushed further sum to a positive value, which the Coulomb law makes zero, and further electrons farther away than distinct ions do not reverse it. At the center of an octahedron, a cube, or a cuboctahedron of ions an electron is pushed off in every direction, and a positive ion is restored in every direction. The ions are held, and the electrons are free. Any cubical truncation of a like-ion simple, body-centered, or face-centered lattice holds the ion and pushes an electron off, to any number of shells, while every phase lies below \(\pi\) and off the first node, and the rate reverses at radius \(\ell/(2\pi)\). The nearest two shells of rock salt, cesium chloride, zincblende, and both sites of fluorite push a positive ion off in the outer well for the ordinary charges, and a crystal site has no fivefold axis.

The square is a local minimum, and its electrons sit on opposite vertices. The gravitational channel supplies the complementary stability: two electron masses driven inside the first gravitational node meet a simple pole, so a shared seat is expelled. The factor \(G m_e^{2}/(k e^{2})=2.401\times 10^{-43}\) thins that neighborhood and does not remove it. The node still prevents the outer-well integral from assigning a dissociation energy, and the Bohr spectrum is not implied. The side of the square is not the Bohr radius. Placing the side at that radius would put the first node strictly between \(91/40\) and \(12/5\) Bohr radii. The laboratory sign at that radius is attraction, so the radius lies in the outer well. That sign forces \(\ell<(19/10)a_0\), and the same sign at \(a_0/Z\) forces \(\ell<(19/10)a_0/Z\). For the input charges \(1\), \(2\), and \(92\) the caps lie strictly between \(1005\times 10^{-13}\) and \(1006\times 10^{-13}\) metres, between \(502\times 10^{-13}\) and \(503\times 10^{-13}\) metres, and between \(1092\times 10^{-15}\) and \(1093\times 10^{-15}\) metres. Placing the first node at \(a_0\) puts the radius \(a_0/2\) where \(\gamma_s<0\). The electromagnetic length is \(\hbar/(m_e c)\), strictly below the charge-\(92\) cap, between \(3861\times 10^{-16}\) and \(3862\times 10^{-16}\) metres, with its first node between \(2032\times 10^{-16}\) and \(2069\times 10^{-16}\) metres and strictly inside \(a_0/92\). The square lies in the reversed shell. The circular radius, the first node, the diagonal, the side, the one-electron extremum, and the second node stand in that order, each outside the next.

A restoring radius is a sign boundary. The radial oscillation across it has \(\omega^{2}=K k e^{2}/(m_e\ell^{3})\), with \(K=4.556\) on the charge-\(8\) dodecahedron. On the square the ratio of the electron wavelength to the diagonal wavelength is \(45.08\), the square root of the ratio of \(K/m\) on the electron coordinate to \(K/m\) on the diagonal, and the ratio does not require an identification of the side. The heat ray of one restoring root is the acceleration of the charge at the period that root already has. For a spacing \(x\ge 1\) the mean occupation lies below \(2e^{-x}\). A transition frequency is an eigenvalue difference and does not scale with temperature. A stationary orbital fixes the electron's motion, so the velocity of the electron relative to the proton is not a quantity the orbital computes, and that difference does not change with the temperature. The continuous spectrum is the occupation-weighted integral over an enclosure whose modes grow as the square of the frequency. Its wavelength density has one maximum between \(9/2\) and \(5\) in units of \(k_B T\). The peak of a body at \(300\,\mathrm{K}\) lies at \(0.1283\,\mathrm{eV}\), the laboratory value of that maximum, and the peak wavelength times the temperature is constant. The integrated frequency density scales as \(T^{4}\), while one classical root scales as \(T\). The mode count and \(h\) are inputs.

\begin{thebibliography}{9}

\bibitem{dst-main}
Anonymous,
``Gravity as Torsion between Dual Spacetime: A Biquaternionic Reformulation of General Relativity'' (2025).
\url{https://github.com/hypernumbernet/dual-spacetime-doc/blob/main/double-spacetime-theory.pdf}.

\bibitem{codata2018}
E.~Tiesinga, P.~J.~Mohr, D.~B.~Newell, and B.~N.~Taylor,
``CODATA recommended values of the fundamental physical constants: 2018,''
\textit{Rev.\ Mod.\ Phys.} \textbf{93}, 025010 (2021).

\end{thebibliography}

\end{document}
